\chapter{Conclusión}
A raíz de lo anteriormente observado, se puede concluir que la implementación de los algoritmos de ordenamiento iterativos en el sistema de gestión de datos de deportistas permitió no solo ordenar efectivamente la información, sino también realizar un análisis comparativo entre diferentes algoritmos. A través de casos de prueba se pudo evaluar el rendimiento de cada algoritmo en distintos escenarios, lo que proporcionó una comprensión más profunda de sus características y comportamientos.

Los resultados obtenidos en los benchmarks mostraron que aunque algunos algoritmos pueden tener un rendimiento similar en ciertos casos, otros pueden presentar diferencias significativas dependiendo del tamaño de la entrada y la naturaleza de los datos. Esto resalta la importancia de elegir el algoritmo adecuado para cada situación específica.

\chapter{Anexos}
\section{Código fuente}
El código fuente del programa desarrollado se encuentra disponible publicamente en el siguiente repositorio de GitHub: \url{https://github.com/ivnmansi/algorithm_analyzer}.
\section{Contribución por integrante}
\begin{itemize}
    \item \textbf{Franco Aguilar}: Elementos base del programa (estructura de datos, gnuplot, lógica de menús, diseño general del sistema, algoritmo de mezcla fisher-yates), revisión de documentación, creación de casos de prueba para el análisis de datos experimentales, implementación de insertion sort y búsqueda secuencial.
    \item \textbf{Iván Mansilla}: Elementos base del programa (estructura de datos, printing, parseo de CSV, comandos de terminal, base de gnuplot, makefile, diseños general del sistema) e implementación de Bubble Sort optimizado y Selection Sort optimizado.
    \item \textbf{Ayrton Morrison}: Elementos base del programa (estructuras de datos, generación y carga de datos, lógica de menús, diseño general del sistema, benchmarks), implementación de cocktail shaker sort y búsqueda binaria.
\end{itemize}